\begin{helper}[A point lies in the interior of at most one edge]
  Let $k \in \mathbb{N}$, let $s \colon \mathbb{N} \to \mathbb{R}$ be monotone (non-decreasing) on
  $\set{0,1,\dots,k}$, and let $c \in \mathbb{R}$. Then
  \[
    \#\set{ j \in \set{1,\dots,k} \colon s(j-1) < c < s(j) } \le 1
    .
  \]
  (Here $j-1$ denotes truncated natural-number subtraction; since $j \ge 1$ in the index set, this
  is ordinary subtraction throughout.)

  \paragraph{Role.} This is the ``edges have pairwise disjoint interiors'' principle in the form
  needed to count the edges that a window's endpoint can cut: a single real number can lie
  strictly inside at most one of the edges $[s(j-1),s(j)]$. It is a statement about the breakpoint
  sequence alone, with no reference to any ambient curve.
\end{helper}
\begin{proof}
  First we record the form of monotonicity we use: for all $p, q \in \mathbb{N}$ with $p \le q$
  and $q \le k$, both $p$ and $q$ lie in $\set{0,\dots,k}$ (the lower bound being automatic in
  $\mathbb{N}$, and $p \le q \le k$), so monotonicity of $s$ on $\set{0,\dots,k}$ gives
  $s(p) \le s(q)$.

  Write $S := \set{ j \in \set{1,\dots,k} \colon s(j-1) < c < s(j) }$. It suffices to show that any
  two elements of $S$ are equal. So let $x, y \in S$ and suppose, for contradiction, that
  $x \ne y$. By trichotomy either $x < y$ or $y < x$; the two cases are symmetric, so assume
  $x < y$ (otherwise swap the names of $x$ and $y$).

  From $x \in S$ we have $1 \le x \le k$ and $s(x-1) < c < s(x)$; from $y \in S$ we have
  $1 \le y \le k$ and $s(y-1) < c < s(y)$.

  Since $x < y$ we have $x \le y-1$, and $y - 1 \le y \le k$, so the recorded monotonicity gives
  \[
    s(x) \le s(y-1)
    .
  \]
  Combining with the two strict inequalities $c < s(x)$ (from $x \in S$) and $s(y-1) < c$ (from
  $y \in S$) yields
  \[
    c < s(x) \le s(y-1) < c
    ,
  \]
  i.e.\ $c < c$, a contradiction. Hence no two distinct elements of $S$ exist, and
  $\#S \le 1$.
\end{proof}
